\begin{claim}{D.30}{Refuted}{A closed genus-two counterexample to full profinite recognition}
The assertion that every torsion-free finite extension of a full
profinite closed hyperbolic surface group is a completed surface group
is false.  There is an exact extension
\[
 1\longrightarrow H=\widehat\pi_1(\Sigma_2)
 \longrightarrow E\longrightarrow C_2\longrightarrow1
                                                        \tag{D.30.1}
\]
with $E$ torsion-free and with a faithful outer action, such that
\[
 H^2(E;\mathbb F_3)=\mathbb F_3,\qquad
 H^2(E;\mathbb F_5)=0.                                  \tag{D.30.2}
\]
No profinite completion of a discrete surface group has these two
cohomology groups.  The exact maximal pro-$p$ quotients are
\[
 E^{(2)}\cong\widehat{\pi_1(N_3)}^{(2)},\qquad
 E^{(3)}\cong\mathbb Z_3^2,\qquad
 E^{(p)}\cong\widehat F_2^{(p)}\quad(p\ne2,3),           \tag{D.30.3}
\]
where $N_3$ is the closed nonorientable surface of genus three.
Thus the counterexample also gives the precise maps through which
its separate pro-$p$ surface presentations arise.
Its outer involution is conjugate to no element of the extended
geometric completion. Its fixed discretization and
prediscretization sets are both empty.
\end{claim}
\begin{proof}
We use the commutator convention $[x,y]=xyx^{-1}y^{-1}$.
Set
\[
 \epsilon_3=-1,\quad\epsilon_p=1\ (p\ne3),\qquad
 \kappa_2=1,\quad\kappa_p=0\ (p\ne2),\qquad
 \lambda=-\epsilon
                                                        \tag{D.30.4}
\]
in $\widehat{\mathbb Z}=\prod_p\mathbb Z_p$.  These elements
satisfy $\epsilon^2=\lambda^2=1$ and
$\epsilon\kappa=\kappa$.

\emph{An actual punctured-torus isomorphism.}
Take the elliptic curve $Y^2=4X^3-4X$ over $\mathbb Q$,
with its origin removed and its rational Weierstrass tangential
base point.  Nakamura's full profinite construction
\cite[\S1, equation (1.2), and \S2.5, p.~424]{C_Nakamura2013}
gives an actual homomorphism
\[
 G_{\mathbb Q}\longrightarrow\operatorname{Aut}
 \langle x,y,z\mid[x,y]z=1\rangle^{\widehat{\ }},
 \qquad \sigma(z)=z^{\chi(\sigma)}.
                                                        \tag{D.30.5}
\]
Here $\chi$ is the full cyclotomic character.  The local
coordinate at the rational puncture and its Kummer covers identify
its inertia with $\widehat{\mathbb Z}(1)$, which explains the exact
power equation.  Every finite cyclotomic extension has Galois group
$(\mathbb Z/n)^\times$; the resulting surjections and compactness
show that $\chi:G_{\mathbb Q}\to\widehat{\mathbb Z}^{\times}$
is onto.  Choose $\sigma$ with $\chi(\sigma)=\lambda$.
Renaming the target generators gives a continuous isomorphism
\[
 \tau:F_1=\widehat F(x,y)\xrightarrow{\sim}
       F_2=\widehat F(u,v),
 \qquad \tau([x,y])=[u,v]^\lambda.                     \tag{D.30.6}
\]
This uses a single Galois element of prescribed multiplier.
The construction below proves its outer involution directly from
$\tau$ and $\tau^{-1}$.

\emph{The automorphism of the closed surface group.}
Use its defining profinite presentation
\[
 H=\langle x,y,u,v\mid[x,y][u,v]=1\rangle^{\widehat{\ }},
 \qquad l=[x,y]^{-1}=[u,v],\qquad c=l^\kappa.           \tag{D.30.7}
\]
Define $\alpha$ on the first pair by the two words of $\tau$,
and on the second pair by the two words of
$\operatorname{Inn}(c)\tau^{-1}$.
Equation (D.30.6) and $\lambda^{-1}=\lambda$ give
\[
 \alpha([x,y])=l^\lambda,\qquad
 \alpha([u,v])=c\,[x,y]^\lambda c^{-1}=l^{-\lambda}.
\]
Their product is one, so this substitution defines a continuous
endomorphism of the presented group $H$.  It gives
$\alpha(l)=l^\epsilon$ and $\alpha(c)=c$.
Substitution on each of the four generators, using the actual inverse
$\tau^{-1}$, then gives
\[
                     \alpha^2=\operatorname{Inn}(c).
                                                        \tag{D.30.8}
\]
Consequently $\operatorname{Inn}(c^{-1})\alpha$ is its inverse.
These identities use the universal property of the displayed
presentation, so no assertion about efficiency of a profinite
amalgam is required.

Define $E$ on the compact space $H\times\{0,1\}$ by
\[
 (h,i)(k,j)=
 \bigl(h\alpha^i(k)c^{\lfloor(i+j)/2\rfloor},
                           i+j\bmod2\bigr).            \tag{D.30.9}
\]
For quotient labels $(1,1,0)$ associativity is precisely
$\alpha^2(k)c=ck$; for $(1,1,1)$ it also uses
$\alpha(c)=c$.  The other six label triples follow by direct
substitution.  The identity is $(1,0)$; the inverse of $(h,0)$ is
$(h^{-1},0)$, and that of $(h,1)$ is
$(\alpha^{-1}(h^{-1}c^{-1}),1)$.
All operations are continuous.  Thus $E$ is a profinite group
and its projection gives (D.30.1).  Write $t=(1,1)$, so that
$tht^{-1}=\alpha(h)$ and $t^2=c$.

\emph{An exact presentation and torsion-freeness.}
Set $u'=\tau(x)$ and $v'=\tau(y)$.  This is a free basis
of the formal second punctured-torus group, with
\[
 [u',v']=[u,v]^\lambda=l^\lambda,\qquad
       u'=txt^{-1},\quad v'=tyt^{-1}.
\]
The relations for $t$ on $u',v'$ follow from $t^2=c$.
Eliminating these two generators therefore gives the exact
profinite presentation
\[
 E\cong\langle x,y,t\mid
          tlt^{-1}=l^\epsilon,\quad t^2=l^\kappa
                 \rangle^{\widehat{\ }},
                 \qquad l=[x,y]^{-1}.                 \tag{D.30.10}
\]
For completeness, the homomorphism to (D.30.9) sends $x,y,t$
to the named elements and is onto because their conjugates
$u',v'$ generate the second factor.  In the reverse direction,
recover $u,v$ by applying the words of $\tau^{-1}$ to
$txt^{-1},tyt^{-1}$.  Their commutator is
$[txt^{-1},tyt^{-1}]^\lambda=l$, by (D.30.10).
Thus the surface relator holds, and the conjugation relations on
both pairs follow from $t^2=c$ and the defining words of
$\alpha$.  The two maps fix all named generators and are inverse.

The universal property of the maximal pro-$2$ quotient substitutes
$\epsilon_2=\kappa_2=1$ into (D.30.10).  Its first relation follows
from the second, giving
\[
 E^{(2)}\cong
       \langle x,y,t\mid[x,y]t^2=1\rangle_{\mathrm{pro}\text{-}2}.
                                                        \tag{D.30.11}
\]
This is the rank-three Demushkin presentation with $q=2$ and
$f=\infty$: cyclically move $t^2$ to the beginning and replace
$x,y$ by their inverses to match the commutator convention of
\cite[\S1.4 and Theorem~1.9]{R-Wilkes}.  The converse in that
theorem proves the Demushkin assertion.  Its abelianization is
$\mathbb Z_2^2\times C_2$, so it is infinite.  The same source's
preceding cohomological-dimension statement gives
$\operatorname{cd}_2(E^{(2)})=2$.  A subgroup of order two has
infinite $2$-cohomological dimension, so $E^{(2)}$ is torsion-free.

There is also an explicit free-basis identification with the
three-crosscap word.  The substitution
\[
 x\mapsto a d^{-1}b^{-1}a^{-1},\qquad
 y\mapsto a b d a b a^{-1},\qquad t\mapsto abd
                                                        \tag{D.30.12}
\]
sends $[x,y]t^2$ to $a^2b^2d^2$.  Its inverse is
\[
 a\mapsto xt,\qquad
 b\mapsto t^{-1}x^{-1}t^{-1}yxt,\qquad
 d\mapsto t^{-1}x^{-1}y^{-1}t^2.
\]
Free reduction checks the relator and both composites.  The maps
extend to mutually inverse pro-$2$ isomorphisms, proving the first
comparison in (D.30.3).

The parity quotient $E\to C_2$ factors through $E^{(2)}$ and
sends $t$ to its nonidentity element and $x,y$ to the identity.
Every finite-order element of $E$ maps trivially to the torsion-free
$E^{(2)}$, hence has trivial parity and lies in $H$.
The latter is torsion-free by A.7.  This proves torsion-freeness of
$E$ itself and, in particular, the absence of a section of (D.30.1).

For an odd prime $p$, (D.30.10) gives $t^2=1$ in $E^{(p)}$,
so $t=1$.  At $p=3$ the remaining relation is $l=l^{-1}$,
which forces $l=1$, and the quotient is
$\langle x,y\mid[x,y]=1\rangle_{\mathrm{pro}\text{-}3}
=\mathbb Z_3^2$.  At every other odd prime the remaining relation
is the identity $l=l$.  The universal properties in both directions
give the free pro-$p$ group on $x,y$.  This proves all of (D.30.3).

\emph{The orientation line and the obstruction to every surface.}
Let $A\in\operatorname{GL}_2(\widehat{\mathbb Z})$ be the
abelianization matrix of $\tau$.  In the finite Heisenberg group
modulo $n$, the commutator of two elements is the central generator
raised to the determinant of their two abelianization columns.
Applying (D.30.6) there gives $\det A=\lambda\pmod n$ for every
$n$, hence $\det A=\lambda$ in $\widehat{\mathbb Z}$.
On the basis $([x],[y],[u],[v])$ of $H^{\mathrm{ab}}$, inner
conjugation disappears and the matrix of $\alpha$ is
\[
 M=\begin{pmatrix}0&A^{-1}\\A&0\end{pmatrix}.
\]
For $J=\left(\begin{smallmatrix}0&1\\-1&0\end{smallmatrix}\right)$,
the identity $A^tJA=(\det A)J$ gives
\[
 M^t\operatorname{diag}(J,J)M
       =\lambda\operatorname{diag}(J,J),                \tag{D.30.13}
\]
because $\lambda^{-1}=\lambda$.  The companion identity
$M\operatorname{diag}(J,J)M^t=\lambda\operatorname{diag}(J,J)$
also holds by $AJA^t=(\det A)J$.  Since the matrix of $\alpha^*$
on degree-one cohomology is $M^t$, this last identity gives the
multiplier of its cup pairing with all dual maps specified.
The completed surface cup-product comparison of A.2 is natural
for automorphisms.  Its nondegenerate degree-one pairing therefore
shows that $\alpha^*$ acts by $\lambda_p$ on the one-dimensional
$H^2(H;\mathbb F_p)$.  The matrix $M$ exchanges two nonzero direct
summands, so $\alpha$ is not inner.  Equation (D.30.8) thus gives
a faithful outer $C_2$ action.

The orientation-line character is a continuous homomorphism
\[
 \operatorname{Out}(H)\longrightarrow
                    \widehat{\mathbb Z}^{\times}.
\]
Its restriction to the extended geometric mapping class group
has image the two diagonal elements $+1,-1$, and continuity gives
the same bound for its closure. Conjugation in $\operatorname{Out}(H)$
preserves this character because its target is abelian. Our element
has value $\lambda$, with $\lambda_3=1$ and $\lambda_5=-1$, so no
conjugate lies in that closure. D.16's exact transporter formulas
therefore give empty fixed discretization and prediscretization
sets for this actual involution.

For odd $p$, averaging over $C_2$ makes all its higher cohomology
with $\mathbb F_p$-coefficients vanish.  The continuous
Hochschild--Serre spectral sequence of (D.30.1) consequently gives
\[
 H^2(E;\mathbb F_p)\xrightarrow[\mathrm{restriction}]{\sim}
                      H^2(H;\mathbb F_p)^{C_2}.
                                                        \tag{D.30.14}
\]
At $p=3$ the scalar is $\lambda_3=1$, while at $p=5$ it is
$\lambda_5=-1$.  This proves (D.30.2).

Finally, the completion of a closed orientable surface group of
positive genus has one-dimensional $H^2(-;\mathbb F_p)$ at every
prime, by the surface comparison of A.2 (including its torus
calculation).  A closed nonorientable surface of genus at least
two has an orientable double cover.  Exact completion at index
two and the same restriction map (D.30.14) compute its degree-two
cohomology at odd primes as the invariants of the orientation
reversal, hence zero.  The sphere group is trivial and the
projective-plane group is $C_2$, so both have zero degree-two
cohomology at odd primes.  A surface with boundary or a noncompact
surface has a free fundamental group, whose profinite completion
has cohomological dimension at most one.  These cases exhaust the
surface groups and exclude the pair of values in (D.30.2).

The independently executed program
\path{code/part_b/closed_extension_d30_audit.py}
checks the free-word isomorphism in (D.30.12), the intersection
matrices and the finite cochain computation of the two invariant
lines.  The profinite construction, presentation comparisons and
torsion-freeness are proved above; the finite computations verify
the stated identities.
\end{proof}
